\documentclass[10pt,a4paper]{amsart}
\usepackage[margin=19mm]{geometry}
\usepackage{amsmath,amssymb,mathtools}
\usepackage[scr=rsfs,cal=euler]{mathalfa}
\usepackage{xfrac}
\usepackage[unicode,hidelinks,bookmarks=false]{hyperref}
\newcommand{\ie}{i.e.\ }
\newcommand{\cf}{cf.\ }
\newcommand{\resp}{respectively\ }
\newcommand{\Subsection}{\subsection}
\providecommand{\nobreakdash}{\nobreak\hbox{-}}
\setlength{\emergencystretch}{2em}
\multlinegap0pt
\usepackage[shortlabels]{enumitem}
\usepackage{amssymb}
\newtheorem{proposition}{Proposition}
\newcommand{\Hh}{\mathbb H}
\title{Positive \texorpdfstring{$\tau$}{tau}-bi-Ricci curvature and Mean curvature flow with surgery in hyperbolic space}
\author{Tianci Luo, Yong Wei, Rong Zhou}
\date{Source: arXiv:2608.25971v1 (CC BY 4.0)}
\begin{document}
\maketitle

\noindent\emph{Source and licence.} Positive \texorpdfstring{$\tau$}{tau}-bi-Ricci curvature and Mean curvature flow with surgery in hyperbolic space. Version arXiv:2608.25971v1 (CC BY 4.0), distributed by arXiv under the Creative Commons Attribution 4.0 International licence, http://creativecommons.org/licenses/by/4.0/, as recorded in the arXiv licence metadata for this exact version and shown on the abstract page https://arxiv.org/abs/2608.25971v1. Full licence notice: \texttt{licenses/arxiv-2608.25971-CC-BY-4.0.txt}.\par\smallskip

\noindent\textit{Editorial scope.} Counted unit: the article's proposition prop:lp-bound (source0.tex lines 1480--1495). The article's complete original proof of it is delivered with it (source0.tex lines 1497--1611, one \texttt{proof} environment of 115 lines). No mathematics is written, repaired or re-derived: the statement and the proof are the article's own lines, in the article's own notation, and no retained line is substituted or re-flowed.  Retained uncounted as the source-local context the proof needs: lines 391--409; lines 504--508; lines 879--882; lines 1088--1094; lines 1162--1164; lines 1260--1282; lines 1408--1413; lines 1422--1434.  Nothing was omitted from any retained span, so the package carries no omission record.  No retained line contains a citation, so the wrapper carries no bibliography and no bibliography entry.  No retained line contains a figure environment, an image-inclusion command or drawing code, so no image is delivered and the package declares no asset dependency.  Editorial formatting only, in addition: an AMS \texttt{amsart} preamble that restates the article's own declarations and macros used by the retained lines, namely the environments \texttt{proposition} on separate counters, together with the standard packages those lines need.  The retained environments print in this document as Proposition 1 in order of appearance, whereas the article prints the same items with the numbering of its own full document.  The complete native source of the article as distributed in the arXiv e-print for this exact version is bundled unchanged as \texttt{sources/208609/source0.tex}.
\begin{align}
 \partial_tg_{ij}=&~-2Hh_{ij},\nonumber\\
 \partial_t d\mu=&~-H^2d\mu, \label{s2.evl-dmu}\\
    \partial_th_{ij}
 =&~\nabla_i\nabla_jH-H\bigl((h^2)_{ij}+\kappa^2g_{ij}\bigr)\nonumber\\
 \partial_th_{ij}  =&~\Delta h_{ij}-2H\bigl((h^2)_{ij}+\kappa^2g_{ij}\bigr)
  +(|A|^2+n\kappa^2)h_{ij},\label{eq:h-direct-evolution}\\
 (\partial_t-\Delta)H=&~(|A|^2-n\kappa^2)H,\label{eq:H-evol}\\
 (\partial_t-\Delta)H^2=&-2|\nabla H|^2
   +2H^2(|A|^2-n\kappa^2),\label{eq:H2-evol}\\
 (\partial_t-\Delta)|A|^2=&-2|\nabla A|^2
   +2|A|^2(|A|^2-n\kappa^2)
   +4n\kappa^2\left(|A|^2-\frac1nH^2\right)\label{eq:A2-evol}\\
 (\partial_t-\Delta)|\nabla A|^2
 \leq&-2|\nabla^2A|^2+c(|A|^2+\kappa^2)|\nabla A|^2,\label{eq:derivative-evolutions}\\
 (\partial_t-\Delta)|\nabla^2A|^2
 \leq&-2|\nabla^3A|^2
 +c\bigl((|A|^2+\kappa^2)|\nabla^2A|^2+|\nabla A|^4\bigr), \label{eq:derivative-evolutions-2}
\end{align}

\begin{equation}\label{eq:uniform-taubiric}
 P_{ij}^{(\tau)}
 \geq\alpha_0 H|\lambda_p-\lambda_q|
 \qquad(i\ne j,\ p,q\text{ arbitrary}).
\end{equation}

\begin{equation}\label{eq:kato-hs}
 |H\nabla_i h_{kl}-(\nabla_iH)h_{kl}|^2
 \geq\frac{\beta^2}{8}H^2|\nabla H|^2.
\end{equation}

\begin{equation}\label{eq:normalized-space-form-setting}
 R=1,\qquad
 0\leq\kappa\leq1,\qquad
 H\geq n\kappa+\alpha_2,\qquad
 |A_0|^2\leq1,\qquad
 |M_0|\leq\alpha_1.
\end{equation}

\begin{equation}\label{eq:acylindrical-region}
 |A|^2-\frac1{n-1}H^2\geq\eta H^2.
\end{equation}

\begin{proposition}[$\tau$-bi-Ricci Poincar\'e inequality]
\label{prop:tau-poincare}
For every $\eta>0$ there exists
$\gamma=\gamma(n,\alpha_0,\eta)>0$ with the following property.
Let
\begin{equation*}
 X:M^n\longrightarrow\Hh^{n+1}_{\kappa}
\end{equation*}
be a smooth closed two-sided hypersurface immersion satisfying
$H>0$, $H\geq n\kappa$, and \eqref{eq:uniform-taubiric}.  Suppose that
$u\in W^{1,2}(M)$ vanishes almost everywhere outside the set
\eqref{eq:acylindrical-region}.  Then, for every $\varrho\geq1$,
\begin{equation}\label{eq:tau-poincare}
 \gamma\int_Mu^2H^2\,d\mu
 \leq
 \int_M\left(
 \varrho^{-1}|\nabla u|^2
 +\varrho u^2\frac{|\nabla A|^2}{H^2}
 \right)d\mu.
\end{equation}
The constant $\gamma$ is uniform for
$0\leq\tau\leq2$ and $0\leq\kappa\leq1$.
\end{proposition}

\begin{equation}\label{eq:D-controls-gradient}
 \frac{|\nabla A|^2}{H^2}
 \leq 2\frac{|\mathcal D|^2}{H^4}
 +2\frac{|A|^2}{H^4}|\nabla H|^2
 \leq C(n,\alpha_0)\frac{|\mathcal D|^2}{H^4}.
\end{equation}

\begin{align}
 (\partial_t-\Delta)f_{\sigma,\eta}
 &=2(1-\sigma)\left\langle
       \nabla f_{\sigma,\eta},\frac{\nabla H}{H}\right\rangle
 -\frac{2}{H^{4-\sigma}}
  |H\nabla_i h_{kl}-(\nabla_iH)h_{kl}|^2 \notag\\
 &\quad-\sigma(1-\sigma)f_{\sigma,\eta}
          \frac{|\nabla H|^2}{H^2}
 +\sigma f_{\sigma,\eta}(|A|^2-n\kappa^2)\notag\\
 &\quad+4n\kappa^2H^{\sigma-2}
       \left(|A|^2-\frac1nH^2\right).
 \label{eq:f-evol}
\end{align}

\begin{proposition}[Uniform $L^p$ bound]\label{prop:lp-bound}
For every $\eta>0$ there are constants $c_4,c_5>0$, depending only on $n$,
$\alpha_0$, and $\eta$, such that if
\begin{equation*}
 p\geq c_4,
 \qquad
 0<\sigma\leq\min\left\{\frac12,c_5p^{-1/2}\right\},
\end{equation*}
then, for $0\leq t<T$,
\begin{equation}\label{eq:lp-bound}
 \int_{M_t}f_+^p\,d\mu
 \leq\int_{M_0}f_+^p\,d\mu
 +C(n,\alpha,\eta,\sigma,p)|M_0|t.
\end{equation}
The constant is independent of $\kappa$ and $T$.
\end{proposition}

\begin{proof}
We work on the set $\{f>0\}$ and use a smooth convex approximation of the
positive part.  Multiplying \eqref{eq:f-evol} by $pf_+^{p-1}$, integrating by
parts, and using \eqref{s2.evl-dmu}, we obtain
\begin{align*}
 \frac{d}{dt}\int f_+^p d\mu
 &\leq-p(p-1)\int f_+^{p-2}|\nabla f|^2d\mu
 +2p(1-\sigma)\int f_+^{p-1}
     \left\langle\nabla f,\frac{\nabla H}{H}\right\rangle d\mu\\
 &\quad-2p\int f_+^{p-1}H^{\sigma-4}|\mathcal D|^2d\mu
 +p\sigma\int f_+^p(|A|^2-n\kappa^2)d\mu\\
 &\quad+Cp\kappa^2\int H^\sigma f_+^{p-1}d\mu.
\end{align*}
The nonpositive term
$-p\sigma(1-\sigma)\int f_+^p|\nabla H|^2/H^2$ and the contribution
$-\int H^2f_+^p$ from the evolving measure have only been discarded.

By \eqref{eq:kato-hs},
$|\nabla H|/H\leq C|\mathcal D|/H^2$.  Consequently,
\begin{align*}
 &2p(1-\sigma)\int f_+^{p-1}
   \left|\left\langle\nabla f,\frac{\nabla H}{H}\right\rangle\right|d\mu\\
 \leq&\frac{p(p-1)}4\int f_+^{p-2}|\nabla f|^2d\mu
 +\frac{Cp}{p-1}\int f_+^p\frac{|\mathcal D|^2}{H^4}d\mu\\
 \leq&\frac{p(p-1)}4\int f_+^{p-2}|\nabla f|^2d\mu
 +C\int f_+^{p-1}H^{\sigma-4}|\mathcal D|^2d\mu.
\end{align*}
Here the last step uses $f_+\leq CH^\sigma$.  For all sufficiently large
$p$, the final term is absorbed into the negative $\mathcal D$-term.

The reaction term is bounded by
\begin{equation*}
 p\sigma\int f_+^p(|A|^2-n\kappa^2)
 \leq Cp\sigma\int H^2f_+^p.
\end{equation*}
For the remaining space-form term, Young's inequality gives, for any
$\theta>0$,
\begin{equation*}
 H^\sigma f_+^{p-1}
 \leq\theta H^2f_+^p
 +C_{\theta,p,\sigma}H^{p\sigma-2(p-1)}.
\end{equation*}
The last exponent is nonpositive for $p\geq2$ and
$0<\sigma<1/2$, and $H\geq\alpha_2$ under
\eqref{eq:normalized-space-form-setting}.  Taking $\theta=\sigma$ and using
$0\leq\kappa\leq1$, we find constants $c_0>0$ and $p_0>2$, depending only
on $n$ and $\alpha_0$, such that, for every $p\geq p_0$,
\begin{align}
 \frac{d}{dt}\int_{M_t}f_+^p\,d\mu
 &\leq-c_0p(p-1)\int_{M_t}f_+^{p-2}|\nabla f|^2\,d\mu\notag\\
 &\quad-c_0p\int_{M_t}f_+^{p-1}H^{\sigma-4}
                  |\mathcal D|^2\,d\mu\notag\\
 &\quad+Cp\sigma\int_{M_t}H^2f_+^p\,d\mu
       +C(n,\alpha,p,\sigma)|M_t|.
 \label{eq:first-lp}
\end{align}

To estimate the first positive term in \eqref{eq:first-lp}, we use the Poincar\'e estimate in Proposition~\ref{prop:tau-poincare}. Set $u=f_+^{p/2}$, which vanishes
outside the set \eqref{eq:acylindrical-region}, and apply
Proposition~\ref{prop:tau-poincare} with $\varrho=p^{1/2}$.  Since
\begin{equation*}
 |\nabla u|^2=\frac{p^2}{4}f_+^{p-2}|\nabla f|^2,
\end{equation*}
while \eqref{eq:D-controls-gradient} and $f_+\leq CH^\sigma$ give
\begin{align*}
 \int f_+^p\frac{|\nabla A|^2}{H^2}\,d\mu
 &\leq C\int f_+^p\frac{|\mathcal D|^2}{H^4}\,d\mu\\
 &\leq C\int f_+^{p-1}H^{\sigma-4}|\mathcal D|^2\,d\mu,
\end{align*}
we obtain, for every $p\geq2$,
\begin{align}
 \int_{M_t}H^2f_+^p\,d\mu
 &\leq Cp^{3/2}\int_{M_t}f_+^{p-2}|\nabla f|^2\,d\mu
 +Cp^{1/2}\int_{M_t}f_+^{p-1}H^{\sigma-4}
              |\mathcal D|^2\,d\mu,
 \label{eq:poincare-cylindrical}
\end{align}
where $C=C(n,\alpha_0,\eta)$ is independent of $\kappa$.

Finally, substituting \eqref{eq:poincare-cylindrical} into
\eqref{eq:first-lp}, we obtain
\begin{align*}
 \frac{d}{dt}\int_{M_t}f_+^p\,d\mu
 &\leq
 -\bigl(c_0p(p-1)-C\sigma p^{5/2}\bigr)
 \int_{M_t}f_+^{p-2}|\nabla f|^2\,d\mu\\
 &\quad
 -\bigl(c_0p-C\sigma p^{3/2}\bigr)
 \int_{M_t}f_+^{p-1}H^{\sigma-4}
 |\mathcal D|^2\,d\mu\\
 &\quad
 +C(n,\alpha,\eta,p,\sigma)|M_t|.
\end{align*}
Choose $c_4\geq\max\{p_0,2\}$ sufficiently large and then
$0<c_5\leq1/2$ sufficiently small, depending only on $n$, $\alpha_0$, and
$\eta$, so that
\begin{equation*}
 c_0p(p-1)-C\sigma p^{5/2}\geq0,
 \qquad
 c_0p-C\sigma p^{3/2}\geq0
\end{equation*}
whenever
\begin{equation*}
 p\geq c_4,
 \qquad
 0<\sigma\leq c_5p^{-1/2}.
\end{equation*}
It follows that
\begin{equation*}
 \frac{d}{dt}\int_{M_t}f_+^p\,d\mu
 \leq C(n,\alpha,\eta,p,\sigma)|M_t|.
\end{equation*}
Since the area is nonincreasing along the smooth mean curvature flow,
integration over $[0,t]$ yields  \eqref{eq:lp-bound}.
\end{proof}

\end{document}
